\begin{circuitikz}
    \tikzset{
        Dflipflop falling/.style={flipflop,
        flipflop def={t1=$C$, t3=$D$, t6=$Q$, t4={\ctikztextnot{$Q$}},
        c1=1, n1=1}},
        %
        add async PC/.style={flipflop def={%
            tu={\ctikztextnot{\texttt{CLR}}}}},
    }

    % Placement of digital components
    \draw(0, 0)                 coordinate(start)
                                node[Dflipflop falling, add async PC](D0){0}
    ([xshift=0.5cm] D0.pin 6)   node[Dflipflop falling, add async PC, anchor=pin 1](D1){1}
    ([xshift=0.5cm] D1.pin 6)   node[Dflipflop falling, add async PC, anchor=pin 1](D2){2}
    ([xshift=0.5cm] D2.pin 6)   node[Dflipflop falling, add async PC, anchor=pin 1](D3){3}
    ([yshift=4cm] D1)           node[nand port, number inputs=2, rotate=180](NAND0){}

    % Internal connections
    (D0.pin 6)                  to[short] (D1.pin 1)     
    (D1.pin 6)                  to[short] (D2.pin 1)     
    (D2.pin 6)                  to[short] (D3.pin 1)
    (D0.pin 3)                  to[short] ++(0,-.75)
                                coordinate(btm)
                                to[short] (\tikztostart -| D0.pin 4)
                                to[short] (D0.pin 4)
    (D1.pin 3)                  to[short] (\tikztostart |- btm)
                                to[short] (\tikztostart -| D1.pin 4)
                                to[short] (D1.pin 4)
    (D2.pin 3)                  to[short] (\tikztostart |- btm)
                                to[short] (\tikztostart -| D2.pin 4)
                                to[short] (D2.pin 4)
    (D3.pin 3)                  to[short] (\tikztostart |- btm)
                                to[short] (\tikztostart -| D3.pin 4)
                                to[short] (D3.pin 4)

    % Inputs
    (D0.pin 1)                  to[short] ++(-0.5,0)
                                coordinate(C)
                                node[left](){CLK}
    
    % Outputs
    ($(D0.pin 6)!0.5!(D1.pin 1)$)  
                                to[short,*-] ++(0,0.75)
                                coordinate(Q0)
                                node[left](){$Q_0$}
    ($(D1.pin 6)!0.5!(D2.pin 1)$)
                                to[short,*-] (\tikztostart |- Q0)
                                coordinate(Q1)
                                node[left](){$Q_1$}
    ($(D2.pin 6)!0.5!(D3.pin 1)$)
                                to[short,*-] (\tikztostart |- Q0)
                                coordinate(Q2)
                                node[left](){$Q_2$}
    (D3.pin 6)                  to[short] (\tikztostart |- Q0)
                                coordinate(Q3)
                                node[left](){$Q_3$}
    
    % Truncation circuit
    (D0.up)                     to[short] ++(0,1)
                                coordinate(CLR)
                                to[short] (\tikztostart -| D3.up)
                                to[short] (D3.up)
    (D1.up)                     to[short,-*] (\tikztostart |- CLR)
    (D2.up)                     to[short,-*] (\tikztostart |- CLR)
    (Q1)                        to[short] (\tikztostart |- NAND0.in 1)
                                to[short] (NAND0.in 1)
    (Q3)                        to[short] (\tikztostart |- NAND0.in 2)
                                to[short] (NAND0.in 2)
    (NAND0.out)                 to[short, -*] (\tikztostart |- CLR)

;
\end{circuitikz}